\documentclass{article}

\usepackage{fontsize}
\changefontsize{14}

\usepackage[a4paper, top=1.5cm, bottom=1.5cm, left=2.5cm, right=1.5cm]{geometry}

\usepackage[T2A]{fontenc}   % кодировка шрифтов для кириллицы
\usepackage[utf8]{inputenc} % кодировка исходного файла (UTF-8)
\usepackage[english, russian]{babel} % правила переносов и типографики

\usepackage{romannum} % римские цифры

\usepackage{amsmath} % вставки текста в формулы, нумерация формул

\usepackage{siunitx} %значок градуса

\usepackage{makecell} %переносы в таблице



\begin{document}

\pagenumbering{arabic}


%ТИТУЛЬНИК


\thispagestyle{empty}

\begingroup
\setlength{\parindent}{0pt}

\begin{bfseries}

\vspace*{0cm}
Министерство образования и науки РФ

\vspace{2cm}
Московский государственный технический университет имени\\Н.Э. Баумана

\vspace{2cm}
Факультет «Фундаментальные науки»

\vspace{2cm}
Кафедра «Теоретическая механика» имени профессора\\ Н.Е. Жуковского

\vspace{2cm}

Домашнее задание № 1\\
Кинематика точки. Простейшие движения твердого тела.

\end{bfseries}

\vspace{2cm}

\textbf{Вариант} \hspace{0.5cm} 20

\vspace{0.5cm}

\textbf{Группа} \hspace{0.5cm} СМ11-21Б

\vspace{0.5cm}

\textbf{Студент} \hspace{0.5cm} Усачев Е.М. \hspace{1.5cm} \underline{\hspace{3cm}} 
\hspace{1cm}\underline{\hspace{3cm}}\\
\small\hspace*{8.5cm} подпись \hspace{2.7 cm} дата \normalsize

\vspace{0cm}
\textbf{Преподаватель} \hspace{0.5cm} Белый В.С.


\vfill
\textbf{Москва 2026 год}
\vspace{1cm}

\endgroup

%ТИТУЛЬНИК


\newpage

\begin{center}\large \textbf{Кинематика точки} \end{center}

\noindent Дано:\\
$x = C(1 + 2 \cos(Et))$;\quad $y = D \sin(Et)$;\quad $C = 2$ м;\quad $D = 3$ м;\quad $E = \frac{\pi}{2}$.

Сделать чертеж, на котором следует изобразить траекторию точки
и для моментов времени $t=0$ c, и $t_1=1$ с показать её положения на траектории.
Для момента времени $t_1=1$ с рассчитать и показать на чертеже в соответствующих 
масштабах положение точки, векторы скорости и ускорения, их составляющие, используемые при 
задании движения точки в декартовой и полярной системах координат, а также
естественным способом.


\newpage



\textit{\textsf{Кинематические характеристики движения точки в декартовой
системе координат}}

\bigskip

С учетом заданных значений $C$, $D$ и $E$ уравнения движения точки принимают вид
\begin{equation}x = 2 + 4 \cos\left(\frac{\pi}{2}t\right) \text{;\quad} y = 3 \sin\left(\frac{\pi}{2}t\right)\end{equation}
и одновременно представляют собой уравнения траектории в параметрической форме,
где параметром является время $t$.


Тогда из основного тригонометрического тождества следует, что
\begin{equation}\left(\frac{x-2}{4}\right)^2 + \left(\frac{y}{3}\right)^2 = 1.\end{equation}
Получаем уравнение траектории точки — уравнение эллипса в координатной форме (см. рисунок 1).

По полученным уравнениям (1) и (2) построим в выбранном масштабе траекторию
точки и покажем на ней положения точек $M_0\text{ и }M$ соответственно в начальный 
$t=0\text{ и в текущий }t=1\text{с}$ моменты времени, в который определяются
кинематические характеристики движения (рис. 1).

Декартовы координаты точки $(x,y)$ в указанные моменты времени, согласно уравнениям
(1), соответственно равны:

при $t=0 \text{ с   } \quad x_0 = 2 + 4 \cos 0 = 6\text{ м, } y_0 = 3 \sin 0 = 0\text{ м;} $

при $t=1 \text{ с   } \quad x = 2 + 4 \cos\left(\frac{\pi}{2}\right) = 2 \text{ м, } y = 3 \sin\left(\frac{\pi}{2}\right) = 3 \text{ м.}$

Проекции векторов скорости и ускорения точки на оси $(Ox,Oy)$ декартовой системы
координат определим дифференцированием по времени уравнений движения (1) и для
$t=1\text{ с}$ получим:
$$v_x=\dot{x} = -4 \cdot \frac{\pi}{2} \sin\left(\frac{\pi}{2}t\right) = -6,28 \text{ м/с;}$$
$$v_y=\dot{y} = 3 \cdot \frac{\pi}{2} \cos\left(\frac{\pi}{2}t\right) = 0 \text{ м/с;}$$
$$a_x=\ddot{x} = -4 \left(\frac{\pi}{2}\right)^2 \cos\left(\frac{\pi}{2}t\right) = 0 \text{ м/с$^2$;}$$
$$a_y=\ddot{y} = -3 \left(\frac{\pi}{2}\right)^2 \sin\left(\frac{\pi}{2}t\right) = -7,4 \text{ м/с$^2$.}$$



Модули скорости и ускорения точки:
$$v=\sqrt{v_x^2+v_y^2}=\sqrt{(-6,28)^2+0^2}=6,28 \text{ м/с;}$$
$$a=\sqrt{a_x^2+a_y^2}=\sqrt{0^2+(-7,4)^2}=7,4 \text{ м/с$^2$. }$$
\newpage
Выберем масштабы для векторов скорости и ускорения и построим в точке $M$
траектории при $t=1$ с, эти векторы и их составляющие по осям декартовой
системы координат.

\bigskip

\textit{\textsf{Кинематические характеристики движения точки в
полярной системе координат}}

\bigskip

Совместим полярную ось $OP$ с осью абсцисс $Ox$ декартовой системы 
координат (см. рис. 1), радиальную ось $Or$ проведем через точку $M$ в ее
положении, соответствующем $t=1$с, построим трансверсальную ось $Op$ путем
поворота радиальной оси на $\ang{90}$ против хода часовой стрелки. В полученной
таким образом полярной системе координат угол $\varphi$ и полярный радиус $r$
точки определяются через декартовы координаты и принимают в текущий момент времени
$t=t_1=1$с следующие значения:

$$r=\sqrt{x^2+y^2}=\sqrt{2^2+3^2}=3,61 \text{ м;}$$
$$\varphi=\arctg{\frac{y}{x}}=\arctg{\frac{3}{2}}=56,3^\circ $$

Для определения проекций векторов скорости и ускорения точки на радиальную $v_r\text{, }a_r$
и трансверсальную $v_p\text{, }a_p$ оси полярной системы координат воспользуемся
соответствующими формулами через проекции кинематических параметров на оси декартовой
системы координат:

$$v_r=\frac{xv_x+yv_y}{\sqrt{x^2+y^2}}=\frac{2\cdot (-6,28)+3\cdot 0}{3,61} = -3,48 \text{ м/с;}$$
$$v_p=\frac{xv_y-yv_x}{\sqrt{x^2+y^2}}=\frac{2\cdot 0-3\cdot (-6,28)}{3,61} = 5,22 \text{ м/с;}$$
$$a_r=\frac{xa_x+ya_y}{\sqrt{x^2+y^2}}=\frac{2\cdot 0+3\cdot (-7,4)}{3,61} = -6,15 \text{ м/с$^2$;}$$
$$a_p=\frac{xa_y-ya_x}{\sqrt{x^2+y^2}}=\frac{2\cdot (-7,4)-3\cdot 0}{3,61} = -4,1 \text{ м/с$^2$.}$$

Построим в выбранных масштабах радиальные и трансверсальные составляющие векторов скорости
и ускорения соответственно $\vec{v}_r\text{, }\vec{v}_p\text{; }\vec{a}_r\text{, }\vec{a}_p$
(см. рис. 1).

\bigskip

\textit{\textsf{Кинематические характеристики движения  
при естественном способе задания движения точки}}

\bigskip

Поскольку касательная ось $\tau$ в этом случае совпадает с осью, направленной по скорости,
то проекция вектора скорости на касательную $v_\tau$ равна модулю скорости:
$v_\tau = v = 6,28$ м/с.
\newpage
Определим проекцию ускорения точки на касательную ось $\tau$:

$$a_\tau=\frac{a_xv_x+a_yv_y}{v}=\frac{0\cdot (-6,28) + (-7,4)\cdot 0}{6,28}=0\text{ м/с$^2$.}$$

Найдем проекцию на нормальную ось $n$:

$$a_n=\left|\frac{a_yv_x-a_xv_y}{v}\right|=\left|\frac{(-7,4)\cdot (-6,28)-0\cdot 0}{6,28}\right|=7,4\text{ м/с$^2$.}$$

Построим в выбранных масштабах касательную $\bar{a}_\tau$ и нормальную $\bar{a}_n$
составляющие ускорения точки на чертеже (см. рис. 1).

По найденному значению нормального ускорения $a_n$, определим радиус кривизны 
траектории в данной ее точке $M$:

$$\rho = \frac{v^2}{a_n}=\frac{6,28^2}{7,4}=5,33\text{ м.}$$


Вывод: в результате расчетов были получены значения кинематических характеристик движения точки
в момент времени $t_1=1$ с, которые представлены в таблице 1.

\begin{flushright}
  \small\textit{Таблица 1} \hspace*{0.4cm}
\end{flushright}

\vspace*{-0.9cm}

\begin{center}
  Значения характеристик движения точки в момент времени $t=1$с
\end{center}

\begin{tabular}{|c|c|c|c|c|}
  \hline
  \rule{0pt}{1cm} № & \hspace*{0.5cm} Название величины \hspace*{0.5cm}& \makecell{Обозначение\\ величины} & Значение & \makecell{Единицы\\ измерения} \\[0.5cm]
  \hline
  \rule{0pt}{1cm} 1  &\makecell{Координата точки по\\ оси $Ox$} & $x$ & 2 & м \\[0.5cm]
  \hline
  \rule{0pt}{1cm} 2  &\makecell{Координата точки по\\ оси $Oy$} & $y$ & 3 & м \\[0.5cm]
  \hline
  \rule{0pt}{1cm} 3  &\makecell{Проекция скорости на\\ ось $Ox$} & $v_x$ & -6,28 & м/с \\[0.5cm]
  \hline
  \rule{0pt}{1cm} 4  &\makecell{Проекция скорости на\\ ось $Oy$} & $v_y$ & 0 & м/с \\[0.5cm]
  \hline
  \rule{0pt}{1cm} 5  &\makecell{Проекция ускорения на\\ ось $Ox$} & $a_x$ & 0 & м/с$^2$ \\[0.5cm]
  \hline
  \rule{0pt}{1cm} 6  &\makecell{Проекция ускорения на\\ ось $Oy$} & $a_y$ & -7,4 & м/с$^2$ \\[0.5cm]
  \hline
\end{tabular}

\newpage

\begin{flushright}
  \small\textit{Окончание табл. 1} \hspace*{0.4cm}
\end{flushright}

\begin{tabular}{|c|c|c|c|c|}
  \hline
  \rule{0pt}{1cm} № & \hspace*{0.5cm} Название величины \hspace*{0.5cm}& \makecell{Обозначение\\ величины} & Значение & \makecell{Единицы\\ измерения} \\[0.5cm]
  \hline
  \rule{0pt}{1cm} 7  &\makecell{Модуль вектора скорости} & $v$ & 6,28 & м/с \\[0.5cm]
  \hline
  \rule{0pt}{1cm} 8  &\makecell{Модуль вектора ускорения} & $a$ & 7,4 & м/с$^2$ \\[0.5cm]
  \hline
  \rule{0pt}{1cm} 9  &\makecell{Полярный радиус} & $r$ & 3,61 & м \\[0.5cm]
  \hline
  \rule{0pt}{1cm} 10 &\makecell{Угол в полярных\\координатах} & $\varphi$ & 56,3 & град. \\[0.5cm]
  \hline
  \rule{0pt}{1cm} 11 &\makecell{Проекция скорости на\\ радиальную ось} & $v_r$ & -3,48 & м/с \\[0.5cm]
  \hline
  \rule{0pt}{1cm} 12 &\makecell{Проекция скорости на\\ трансверсальную ось} & $v_p$ & 5,22 & м/с \\[0.5cm]
  \hline
  \rule{0pt}{1cm} 13 &\makecell{Проекция ускорения на\\ радиальную ось} & $a_r$ & -6,15 & м/с$^2$ \\[0.5cm]
  \hline
  \rule{0pt}{1cm} 14 &\makecell{Проекция ускорения на\\ трансверсальную ось} & $a_p$ & -4,1 & м/с$^2$ \\[0.5cm]
  \hline
  \rule{0pt}{1cm} 15 &\makecell{Проекция скорости на\\ касательную ось} & $v_\tau$ & 6,28 & м/с \\[0.5cm]
  \hline
  \rule{0pt}{1cm} 16 &\makecell{Проекция ускорения на\\ касательную ось} & $a_\tau$ & 0 & м/с$^2$ \\[0.5cm]
  \hline
  \rule{0pt}{1cm} 17 &\makecell{Проекция ускорения на\\ нормальную ось} & $a_n$ & 7,4 & м/с$^2$ \\[0.5cm]
  \hline
  \rule{0pt}{1cm} 18 &\makecell{Радиус кривизны\\ траектории} & $\rho$ & 5,33 & м \\[0.5cm]
  \hline
\end{tabular}


\newpage

\begin{center}\large \textbf{Кинематика простейших движений твердого тела} \end{center}

\noindent Дано:\\
$s = d t + f t^2$,\quad $d = -10$ см/с,\quad $f = 30$ см/с$^2$,\quad $r_2 = 10$ см,\quad $r_4 = 5$ см.

\medskip

\noindent Необходимо:

-- выполнить с соблюдением масштаба чертеж механизма;

-- для момента времени $t_1=1$c определить значения угловых скоростей
и угловых ускорений всех пронумерованных тел механизма, совершающих вращательное движение,
и показать их на чертеже круговыми стрелками;

-- для момента времени $t_1=1$c для обозначенных на схемах вариантов задания
точек $A_i$, $i$-х тел найти значения векторов скоростей и ускорений
(с их нормальными и тангенциальными составляющими) и изобразить все указанные векторы
на чертеже с соблюдением принятых масштабов.
\bigskip 
\bigskip 
\bigskip 
\bigskip 
\bigskip
\bigskip
\bigskip
\bigskip
\bigskip
\bigskip
\bigskip
\bigskip
\bigskip
\bigskip
\bigskip 
\bigskip 
\bigskip 
\bigskip 
\bigskip
\bigskip
\bigskip
\bigskip
\bigskip
\bigskip
\bigskip
\bigskip
\bigskip
\bigskip
\bigskip
\bigskip
\bigskip
\bigskip
\begin{center}\text{Рис. 2. Схема механизма} \end{center}

\newpage


\textit{\textsf{Кинематические характеристики движения точек  
и звеньев механизма}}

\bigskip 


Сначала определим виды движения звеньев данного механизма: звено $1$ совершает поступательное
движение, звено $2$ — вращательное вокруг оси $O_{2z}$, звено $3$ (нить) — поступательное, а звено
$4$ — вращательное движение вокруг оси $O_{4z}$ (см. рисунок 3).

Уравнение (3) поступательного движения ведущего звена $1$ с учетом заданных значений
$d$ и $f$ принимает вид:

\begin{equation}
  s = -10t + 30t^2\text{.}
\end{equation}

Дифференцируя по времени уравнение (3), найдем скорость и ускорение звена 1 для момента времени 
$t=1$c:

$$v_1=\dot{s}= -10 + 60t = -10 + 60\cdot 1 = 50 \text{ см/с;}$$

$$a_1=\ddot{s}=60 \text{ см/с$^2$.}$$

\medskip

Точка $A_1$ принадлежит звену 1.

Модуль скорости точки $A_1$:
$$v_{A_1}=v_1=50\text{ см/с.}$$

Модуль ускорения точки $A_1$:
$$a_{A_1}=a_1=60\text{ см/с$^2$.}$$

Звено 2 совершает вращательное движение. Ввиду отсутствия проскальзывания звена 1
по звену 2 скорости и касательные составляющие ускорений точек контакта этих тел равны, поэтому

$$\omega_2=\frac{v_1}{r_2}=\frac{50}{10}=5\text{ рад/с}$$
\noindent--- модуль угловой скорости звена 2;

$$\varepsilon_2=\frac{a_1}{r_2}=\frac{60}{10}=6\text{ рад/с$^2$}$$
\noindent--- модуль углового ускорения звена 2.

\medskip

Точка $A_2$ принадлежит звену 2.

Модуль скорости точки $A_2$:
$$v_{A_2}= \omega_2 r_2 = 5 \cdot 10 = 50\text{ см/с.}$$

Полное ускорение точки $A_2$:
$$\bar{a}_{A_2}=\bar{a}_{A_2}^\tau+\bar{a}_{A_2}^n\text{.}$$

\newpage

Модуль тангеницального ускорения точки $A_2$:
$$\left|\:\bar{a}_{A_2}^\tau\right| = \varepsilon_2 r_2 = 6 \cdot 10 = 60\text{ см/c$^2$.}$$

Модуль нормального ускорения точки $A_2$:
$$a_{A_2}^n=\omega_2^2 r_2 = 5^2 \cdot 10 = 250\text{ см/с$^2$.}$$

Вектор нормального ускорения $\bar{a}_{A_2}^n$ направлен к центру кривизны траектории точки, т.е. к точке $O_2$, расположенной на оси вращения звена 2, а вектор тангенциального ускорения $\bar{a}_{A_2}^\tau$ расположен на касательной к окружности радиуса $r_2$ и соответствует направлению  углового ускорения звена 2. 

Модуль полного ускорения точки $A_2$:
$$a_{A_2}=\sqrt{(a_{A_2}^\tau)^2+(a_{A_2}^n)^2}=\sqrt{60^2+250^2}=257,1\text{ см/с$^2$.}$$

Звено 3 (нить) совершает поступательное движение. Так как нить по условию нерастяжима и движется без проскальзывания, то модули скорости и ускорения звена 3 соответственно равны модулям скорости и тангенциального ускорения точки $A_2$:
$$v_3 = v_{A_2} = 50 \text{ см/с;} \quad a_3 = a_{A_2}^\tau = 60 \text{ см/с$^2$.}$$

Точка $A_3$ принадлежит звену 3.

Модуль скорости точки $A_3$:
$$v_{A_3}=v_3 = 50\text{ см/с.}$$

Модуль ускорения точки $A_3$:
$$a_{A_3}=a_3 = 60\text{ см/с$^2$.}$$

Звено 4 совершает вращательное движение. Так как нить по условию нерастяжима и движется без проскальзывания, то модули касательных скорости и ускорения звена 4 соответственно равны модулям скорости и ускорения звена 3:
$$v_4 = v_3 = 50 \text{ см/с;} \quad a_{A_4}^\tau = a_3 = 60 \text{ см/с$^2$.}$$

Найдем модули угловой скорости и углового ускорения звена 4:
$$ \omega_4 = \frac{v_4}{r_4} = \frac{50}{5} = 10 \text{ рад/с;} \quad \varepsilon_4 = \frac{a_{A_4}^\tau}{r_4} = \frac{60}{5} = 12 \text{ рад/с$^2$.} $$

Точка $A_4$ принадлежит звену 4.

Модуль скорости точки $A_4$:
$$v_{A_4}= \omega_4 r_4 = 10 \cdot 5 = 50 \text{ см/с.}$$

Найдем нормальное и полное ускорение точки $A_4$:
$$ a_{A_4}^n = \omega_4^2 r_4 = 10^2 \cdot 5 = 500 \text{ см/с$^2$;}$$
$$ a_{A_4} = \sqrt{(a_{A_4}^\tau)^2 + (a_{A_4}^n)^2} = \sqrt{60^2 + 500^2} = 503,6 \text{ см/с$^2$.} $$

\newpage

Вектор нормального ускорения $\bar{a}_{A_4}^n$ направлен к центру кривизны траектории точки, т.е. к точке $O_4$, расположенной на оси вращения звена 4, а вектор тангенциального ускорения $\bar{a}_{A_4}^\tau$ расположен на касательной к окружности радиуса $r_4$ и соответствует направлению  углового ускорения звена 4.

Вывод: в результате расчетов были получены значения кинематических характеристик рассматриваемого
механизма в момент времени $t_1=1$ с, которые представлены в таблице 2.

\begin{flushright}
  \small\textit{Таблица 2} \hspace*{0.4cm}
\end{flushright}

\vspace*{-0.9cm}

\begin{center}
  Значения кинематических характеристик в момент $t_1=1$с
\end{center}

\begin{tabular}{|c|c|c|c|c|}
  \hline
  \rule{0pt}{1cm} № & \hspace*{0.5cm} Название величины \hspace*{0.5cm} & \makecell{Обозначение\\ величины} & Значение & \makecell{Единицы\\ измерения} \\[0.5cm]
  \hline

  \rule{0pt}{1cm} 3  & \makecell{Угловая скорость\\звена 2} & $\omega_2$ & 5 & рад/с \\[0.5cm]
  \hline
  \rule{0pt}{1cm} 4  & \makecell{Угловая скорость\\звена 4} & $\omega_4$ & 10 & рад/с \\[0.5cm]
  \hline
  \rule{0pt}{1cm} 5  & \makecell{Угловое ускорение\\звена 2} & $\varepsilon_2$ & 6 & рад/с$^2$ \\[0.5cm]
  \hline
  \rule{0pt}{1cm} 6  & \makecell{Угловое ускорение\\звена 4} & $\varepsilon_4$ & 12 & рад/с$^2$ \\[0.5cm]
  \hline
  \rule{0pt}{1cm} 7  & \makecell{Скорость\\точки $A_1$} & $v_{A_1}$ & 50 & см/с \\[0.5cm]
  \hline
  \rule{0pt}{1cm} 8  & \makecell{Скорость\\точки $A_2$} & $v_{A_2}$ & 50 & см/с \\[0.5cm]
  \hline
  \rule{0pt}{1cm} 9  & \makecell{Скорость\\точки $A_3$} & $v_{A_3}$ & 50 & см/с \\[0.5cm]
  \hline
  \rule{0pt}{1cm} 10 & \makecell{Скорость\\точки $A_4$} & $v_{A_4}$ & 50 & см/с \\[0.5cm]
  \hline
  \rule{0pt}{1cm} 11 & \makecell{Ускорение\\точки $A_1$} & $a_{A_1}$ & 60 & см/с$^2$ \\[0.5cm]
  \hline
  \rule{0pt}{1cm} 20 & \makecell{Тангенциальное\\ускорение точки $A_2$} & $a_{A_2}^\tau$ & 60 & см/с$^2$ \\[0.5cm]
  \hline
  
\end{tabular}

\newpage

\begin{flushright}
  \small\textit{Окончание табл. 2} \hspace*{0.4cm}
\end{flushright}

\begin{tabular}{|c|c|c|c|c|}
  \hline
  \rule{0pt}{1cm} № & \hspace*{0.5cm} Название величины \hspace*{0.5cm} & \makecell{Обозначение\\ величины} & Значение & \makecell{Единицы\\ измерения} \\[0.5cm]
  \hline
  \rule{0pt}{1cm} 16 & \makecell{Нормальное\\ускорение точки $A_2$} & $a_{A_2}^n$ & 250 & см/с$^2$ \\[0.5cm]
  \hline
  \rule{0pt}{1cm} 12 & \makecell{Полное ускорение\\точки $A_2$} & $a_{A_2}$ & 257,1 & см/с$^2$ \\[0.5cm]
  \hline
  \rule{0pt}{1cm} 13 & \makecell{Ускорение\\точки $A_3$} & $a_{A_3}$ & 60 & см/с$^2$ \\[0.5cm]
  \hline
  \rule{0pt}{1cm} 22 & \makecell{Тангенциальное\\ускорение точки $A_4$} & $a_{A_4}^\tau$ & 60 & см/с$^2$ \\[0.5cm]
  \hline
  \rule{0pt}{1cm} 18 & \makecell{Нормальное\\ускорение точки $A_4$} & $a_{A_4}^n$ & 500 & см/с$^2$ \\[0.5cm]
  \hline
  \rule{0pt}{1cm} 14 & \makecell{Полное ускорение\\точки $A_4$} & $a_{A_4}$ & 503,6 & см/с$^2$ \\[0.5cm]
  \hline
\end{tabular}

\end{document}
